\documentclass[oneside,final,14pt]{extreport}
\usepackage[english,russian]{babel}
\usepackage[T2A]{fontenc}
\usepackage{cmap}
\usepackage{graphicx}
\usepackage{amsmath}
\usepackage{amssymb}
\usepackage[a4paper, left=1cm, right=1cm, top=2cm, bottom=2cm]{geometry}
\usepackage{multirow}
\graphicspath{ {images/} }

\title{\textbf{Вопросы ЧГК в КиМ 2026-2027}}
\author{Тараканов Павел, Дитятьева Мария ...}

\begin{document}

\maketitle

\section*{Вопросы ЧГК в КиМ 2026-2027}
\begin{enumerate}
    \item Во вселенной Вархаммер в годы ВКП человечество, в большинстве своём, было атеистами, отрицая существование Богов. В том числе люди так же отказывались верить в приметы. Одним из главных поборников такого подхода был Робаут Гиллиман, примерх одного из 20 легионов космодесанта. Он не раз заявлял, что сделает свой легион наиболее успешным, несмотря на его номер. Какой номер?

    Ответ: 13

    Комментарий: примарх заявлял, что вопреки неудачности числа 13 он сделает свой легион самым успешным.

    \item Действие игр серии Assassin's Creed происходит в тщательно воссозданных исторических декорациях. В апреле 2019 года в мировых СМИ появилась информация, что трёхмерная карта и материалы одной из игр серии (а именно Assassin's Creed: Unity) могут быть официально использованы специалистами. Использованы для чего?

    \textbf{надо намекнуть на париж}

    Ответ: Для реконструкции/восстановления собора Парижской Богоматери после разрушительного пожара.

    \item Одной из тайн, волнующих фанатов вселенной молота войны, является тайна 2 и N-ого легиона космодесанта. Записи о них и их примархах были стерты из всех архивов еще до Ереси Хоруса, однако многие считают что они были уничтоженный Леманом Руссом, прозванным палачём императора, и его волками. Среди фанатов ходит шутка, что N-ый легион был уничтожен по ошибки, потому что дикарь Русс перепутал число 2 и N. Чему равно N?

    Ответ: 11

    Комментарий: римская 2 похожа на арабскую 11.

    \item Художник и дизайнер ЙЕрун ван дер Мост разработал программу, которая с помощью вайб-кОдинга меняет цвета ячеек Excel в реальном времени. Недавно в Таллинском музее в рамках экспозиции, посвящённой Excel, был показан фильм. Назовите этот фильм.

    Ответ: «Матрица». Зачёт: по упоминанию «Матрицы».

    Комментарий. Видим код – говорим «Матрица».

    \item У самцов плодовых мушек дрозофИл синтез феромонов усиливается под действием определённого класса веществ. Наибольший же эффект оказывает ОН, делая самцов ослепительно привлекательными для противоположного пола. Назовите ЕГО максимально точно.

    Ответ: метиловый спирт. Зачёт: метанол, древесный спирт, карбинол, метилгидрат, гидроксид метила, CH3OH.

    Комментарий. А человек от метилового спирта может ослепнуть.

    \item Американский «НаутИлус» оказался ТАМ в 1959 году первым, что вполне логично. Советский «Ленинский комсомол» только в 1970-м. А вот китайская АПЛ попала ТУДА только в 2020-м, причем китайцы к этому не имели никакого отношения. А куда – ТУДА?

    Ответ: на почтовую марку. Зачет: на марку.

    Комментарий. Конечно это НЕ Северный полюс. АПЛ «Наутилус» на Северный полюс попала в 1958-м, а в 1959-м её увековечили на почтовой марке. В СССР путь в филателию занял несколько больше времени. А китайская АПЛ 094 «Цзинь» на марки попала только в 2020-м, и то, не Китая, а Либерии.


    \item Этот святой жил предположительно в V-VII веке, прославился отшельничеством, избавил Фивы от чумы. Его имя по-русски было бы палиндромом, но в России его стали писать не через А, а через О. Вероятно, такое имя носил отец персонажа русских сказок. Напишите это ИМЯ.

    Ответ: ПотАп. Зачет: Патап, ПатАпий.

    Комментарий. Древнеегипетское, точнее коптское, имя Патап (данный Аписом) в России стали писать через О — Потап. Применительно к святому обычно используется в греческом варианте — Патапий, Патапиос. Ну, а Михайло Потапыч Топтыгин хорошо известен.

    \item В одной книге изобретение БЕтти НЕсмит Грэм названо «волшебной маскировочной жидкостью». Это изобретение часто называют так же, как и человека… Какой профессии?

    Ответ: корректор.

    Комментарий. Жидкость маскировала опечатки и ошибки.

    \item ОНИ и великолепная память помогли Николаю Пржевальскому добыть средства на первую экспедицию. Некоторые другие ОНИ после путешествий Пржевальского изменились, причём в лучшую сторону. Назовите ИХ одним словом.

    Ответ: карты.

    Комментарий. Пржевальский был успешным картёжником, отлично запоминал
    комбинации. Добыв игрой необходимую сумму, карты навсегда оставил. Путешествия Пржевальского позволили составить либо уточнить географические карты Восточной и Центральной Азии.

    \item \textbf{УКАЗАНИЕ ВЕДУЩЕМУ: СДЕЛАТЬ АКЦЕНТ НА СЛОВАХ «ПЕРВОГО» И «ЧЕТЫРНАДЦАТОГО»}

    На одной карикатуре осуществлению очередного амбициозного проекта ПЕРВОГО мешает Шрек. ПЕРВЫЙ был ЧЕТЫРНАДЦАТЫМ. Назовите имя ПЕРВОГО.

    Ответ: Пётр. Зачёт: точный ответ.

    Комментарий. Видимо, одного города на болоте императору показалось недостаточно. Петр I был четырнадцатым ребенком Алексея Михайловича.

    \item Пакистанский ОН известен меньше, чем ОН соседней страны, и гораздо меньше, чем ОН американский. Поскольку эта индустрия наиболее развита в Лахоре, пакистанскому ЕМУ, по аналогии с двумя другими, дали неофициальное прозвище Лол [ЛОЛ]… Напишите полностью слово, сокращённое в предыдущем предложении.

    Ответ: Лолливуд.

    Комментарий. ОН — кинематограф. По аналогии с более известными американским Голливудом и индийским Болливудом (названным по городу Бомбею), кинематограф пакистанского Лахора называют Лолливудом. Кстати, на заре пакистанского кинематографа женские роли играли мужчины.

    \item Развивая известное сравнение Земли с НЕЮ, РОберт ЗУбрин заметил, что тогда Солнечная система – это школьный двор. ОНА входит в название изобретения, статья Википедии о котором сопровождена несколькими gif-изображениями . Назовите ЕЁ одним словом.

    Ответ: колыбель.

    Комментарий. РОберт ЗУбрин как-то заметил: «Если Земля – колыбель человечества, то Солнечная система – школьный двор». В статье Википедии «Колыбель НьЮтона» для наглядности показано, как работает система шаров.

    \item Лука Луизелли вместе с группой коллег получили в 2025 году Шнобелевскую премию в номинации «Питание» за исследование того, какое влияние на рацион агам оказало соседство с людьми. Выяснилось, что, в целом совпадая во вкусах с известными сородичами, агамы однозначно предпочитают один «квартет» другому. Назовите оба «квартета».

    Ответ: «4 сыра», «4 сезона».

    Комментарий. Обитающие на пляжах Того агамы отказались от своего привычного рациона из насекомых и прониклись горячей любовью к пицце, как и черепашки-ниндзя, но сыр их привлекает больше, чем ветчина, грибы, овощи и оливки. Итальянская фамилия учёного тоже могла вам помочь.

    \item "Не пнешь — не полетит". "Как пнешь, так и полетит". "Как пнешь, так и получишь". Ответьте двумя словами, для запоминания чего предназначены эти три фразы.

    Ответ: Законов Ньютона.

    Комментарий: Вульгарно, но смысл каждого закона запомнить можно.

    \item Математический анекдот. Математика спрашивают: "Есть ли крылья у слона?" - "Есть, - отвечает математик, но они..." Закончите фразу двумя словами!

    Ответ: "... равны нулю".

    Комментарий: комментарии излишни

    \item Когда в студенческой столовой ленинградского университета появились три варианта качественно различных обедов, то эти обеды в порядке убывания цены стали именоваться студентами соответственно действительными, комплексными и... Как назывался третий, самый дешевый вариант?

    Ответ: Мнимый обед

    Комментарий: комплексное число в математике состоит из мнимой и действительной части

    \item Станислав Ежи Лец спрашивает: «Что ты скажешь на это, физика? В отношениях между людьми трения ведут к … . А к чему ведут трения?

    Ответ: К охлаждению

    Комментарий: в физике при трении происходит выделение тепла и нагревание тел.

    \item В 1912 году в газетах появилась сенсационная фотография: на палубе корабля стоят хорошо одетые господа, а над ними в воздухе завис странный аппарат. Это был один из первых авиационных трюков для прессы. Однако физик, взглянув на фото, сразу сказал бы, что это монтаж. Какая физическая несообразность выдала подделку?

    Ответ: Отсутствие ветра / Прически и одежда.

    Комментарий: Если бы дирижабль или самолет висел прямо над палубой на малой высоте, винты создавали бы мощный поток воздуха, который раздувал бы одежду и волосы людей. На фото люди стояли со спокойными прическами, значит, снимок смонтирован.

    \item Один сингапурский математик назвал своего сына именем, популярным в Англии. Когда родился еще один сын, он назвал младшего китайским именем. Можно сказать, что оба сына вместе представляют собой смешанный экстремум. Как зовут каждого из сыновей?

    Ответ: Макс, Мин

    \item Голландский физик Гендрик Лоренц нашел соотношение между плотностью вещества и его преломляющей способностью. Несколько ранее эту же формулу, но совершенно другим путем получил датский физик. За формулой утвердилось двойное название. Открытие этой формулы является любопытным случаем с точки зрения теории вероятности. Как называется эта формула?

    Ответ: Формула Лоренца-Лоренца

    Комментарий: фамилия датского физика тоже была Лоренц

    \item 1. Встретились два математика. Первый спрашивает второго, есть ли у него дети. "Да", — говорит второй, "У меня трое детей". "А сколько им лет?" — спросил Первый, на что Второй ответил: "Сумма их возрастов равна числу окон в этом доме, а произведение — 36". Первый математик посмотрел на дом, подумал и говорит: "Этой информации недостаточно". "Да, конечно, я совсем забыл: младший сын — рыжий". Сколько лет детям (каждому) второго математика?

    Ответ: 1, 6, 6

    Комментарий: Возможны такие варианты сумм при произведении 36: 1+2+18=21; 2+2+9=13; 2+3+6=11; 3+3+4=10; 1+3+12=16; 1+4+9=14; 1+6+6=13. Так как первый математик знал количество окон в доме, но после раздумий эти условия его не удовлетворили, то следовательно ответ не однозначный. Это реализуется в случаях 2 и 7. В случае 2 младшие дети — близнецы, а по условиям задачи младший ребенок отличается от остальных — он рыжий. Следовательно, это вариант 7.

    \item Вы видите знак, которым обозначается единица измерения электрического сопротивления — Ом. Справа — обозначение Сименса — единицы измерения электрической проводимости, величины, обратной Ому. Эта единица была введена Сименсом в 1860 году, а до этого использовалась другая величина. Мы не спрашиваем вас, как она называлась. Нарисуйте знак, которым она обозначалась.

    Ответ: рисунок с презы

    Комментарий: Раньше применялось название "мо" (англ. mho), представляющее собой прочитанное назад слово "ом" (ohm); обозначалось перевернутой буквой $\Omega$

    \item Это теория, в названии которой есть 2 слова: существительное и прилагательное. Одно из них связано с известным героем русских народных сказок, а другое - с самыми мощными вычислительными машинами на Земле.

    \item В 2021 году человек под ником IlluminatiPirate опубликовал на англоязычном форуме эссе на одну очень интересную конспирологическую теорию. Автор утверждал что начиная приблизительно с 2016—2017 годов основная часть контента в интернете создаётся ботами и системами ИИ ,что «живой» интернет где люди, творят и взаимодействуют друг с другом почти исчез. Вместо него все больше появляется имитации жизни: тексты, сгенерированные нейросетями, одинаковые комментарии, фальшивые аккаунты и активность ботов. Все это создает иллюзию бурной активности, за которой все меньше настоящих людей. Радикальная версия теории утверждает, что это является частью координированных усилий по манипулированию общественным сознанием. Однако исследователи и журналисты считают ее скорее метафорой, чем научно обоснованным утверждением , считая что люди просто боятся неконтролируемое ими изменение интернет-пространства.  Какое название у этой теории?

    Ответ: теория мертвого интернета
    \item
    \item
    \item
    \item
    \item
    \item
    \item
    \item
\end{enumerate}

\end{document}
